\section{Retained data on expanding annuli}
\label{sec:scales}

We first arrange the foundation estimates on a sequence of expanding rescaled
annuli.  Throughout this section, let
\begin{equation}
 s_\ell\longrightarrow0,\qquad
 Z_\ell\ge Z_{\mathrm{bar}}(s_\ell),\qquad
 Z_\ell s_\ell^3\longrightarrow\infty,
 \label{scales:admissible}
\end{equation}
and call such a sequence \emph{admissible}. Choose an arbitrary
normalized neutral ground state $\Psi_\ell$ at each
$Z_\ell$.  The threshold $Z_{\mathrm{bar}}$ includes the fixed barrier
requirements of the preceding section.  At each index we use its observations
and master packets, with the barrier parameters fixed once and for all.
Probabilities and conditional expectations refer to the enlarged law of that
system; no common probability space for different indices is needed.  We omit
a finite initial segment whenever necessary.

\begin{proposition}[Estimates on retained data]
\label{scales:good}
There are deterministic integers $q_\ell\to\infty$ and $j_\ell<J_\ell$,
and events $G_\ell\in\mathcal F_{j_\ell}$ with
$\Prob(G_\ell)\to1$, such that
\begin{equation}
 a_\ell:=\frac{r_{j_\ell}}{s_\ell}
 \in\left[\frac1{2q_\ell},\frac1{q_\ell}\right].
 \label{scales:inner-radius}
\end{equation}
Put $\lambda_\ell=Z_\ell s_\ell^3$ and, on each retained datum, define
\begin{equation}
 \begin{aligned}
 F_\ell(x)&=s_\ell^4H_{j_\ell}(s_\ell x),&
 \sigma_\ell(x)&=s_\ell^6\mu_{j_\ell}(s_\ell x),\\
 U_\ell(x)&=s_\ell^4u_{j_\ell}(s_\ell x),&
 p_\ell(x)&=s_\ell^6p_{j_\ell}(s_\ell x).
 \end{aligned}
 \label{scales:fields}
\end{equation}
These fields satisfy the exact identities
\begin{equation}
 \begin{gathered}
 \sigma_\ell\ge0,\qquad \int_{\R^3}\sigma_\ell\,\dd x=\lambda_\ell,
 \qquad F_\ell=\lambda_\ell K-K*\sigma_\ell,\\
 \Delta F_\ell=-4\pi\lambda_\ell\delta_0+4\pi\sigma_\ell,
 \end{gathered}
 \label{scales:mass-field}
\end{equation}
and, simultaneously on $a_\ell\le|x|\le q_\ell$,
\begin{equation}
 F_\ell(x)\le C_{\mathrm{cap}}|x|^{-4},\qquad
 \bigl|\sigma_\ell(x)-k(F_\ell(x)_+)^{3/2}\bigr|
 \le q_\ell^{-8}|x|^{-6}.
 \label{scales:band}
\end{equation}
The scaled subsolution obeys, in distributions on $\R^3$,
\begin{equation}
 \Delta U_\ell\ge
 -4\pi\lambda_\ell\delta_0
 +4\pi\mathbf1_{\{|x|<a_\ell\}}\sigma_\ell
 -4\pi p_\ell
 +\mathbf1_{\{|x|\ge a_\ell\}}f(U_\ell),
 \label{scales:invariant}
\end{equation}
with
\begin{equation}
 B|x|^{-4}\left(1-\frac{a_\ell}{8|x|}\right)
 \le U_\ell(x)\le C_{\mathrm{inv}}|x|^{-4}
 \qquad (|x|\ge a_\ell),
 \label{scales:barriers}
\end{equation}
and
\begin{equation}
 p_\ell\ge0,\qquad
 \operatorname{supp}p_\ell\subset\{|x|\le a_\ell\},\qquad
 \int_{\R^3}p_\ell\,\dd x\le s_\ell^3.
 \label{scales:error}
\end{equation}
For each finite system, $\sigma_\ell$ is a bounded smooth density,
$p_\ell$ is a bounded density, and both
$F_\ell-\lambda_\ell K$ and $U_\ell-\lambda_\ell K$ extend continuously
and locally Lipschitz through the origin.  All these assertions hold for every
datum in $G_\ell$, so they may be used along arbitrary selections of retained
data.
\end{proposition}

Figure~\ref{fig:scales} distinguishes the physical radius $s_\ell$
from the expanding annulus on which its rescaled fields are compared.
\begin{figure}[!ht]
\centering
\begin{tikzpicture}[x=1cm,y=1cm,>=Latex,font=\small]
 \fill[blue!8] (2.8,.83) rectangle (8.1,1.17);
 \fill[blue!8] (2.8,-1.17) rectangle (8.1,-.83);
 \draw[->] (.6,1)--(9,1);
 \draw[->] (.6,-1)--(9,-1);
 \node[anchor=east] at (.45,1) {physical};
 \node[anchor=east] at (.45,-1) {rescaled};
 \foreach \x in {1.2,2.8,5.3,8.1}{
   \draw (\x,.91)--(\x,1.09);
   \draw (\x,-1.09)--(\x,-.91);
 }
 \node[above] at (1.2,1.1) {$r_0$};
 \node[above] at (2.8,1.1) {$r_{j_\ell}$};
 \node[above] at (5.3,1.1) {$s_\ell$};
 \node[above] at (8.1,1.1) {$q_\ell s_\ell$};
 \node[below] at (1.2,-1.1) {$r_0/s_\ell$};
 \node[below] at (2.8,-1.1) {$a_\ell\asymp q_\ell^{-1}$};
 \node[below] at (5.3,-1.1) {$1$};
 \node[below] at (8.1,-1.1) {$q_\ell$};
 \draw[->,gray] (5.3,.65)--(5.3,-.6)
   node[midway,right,black] {divide radii by $s_\ell$};
\end{tikzpicture}
\caption{The physical and rescaled radii in
Proposition~\ref{scales:good} (schematic, not to scale).
The same observation index supplies both the comparison on the shaded
annulus and the intermediate subsolution.  Its rescaled inner radius
tends to zero while the outer endpoint tends to infinity.}
\label{fig:scales}
\end{figure}


\begin{proof}
\emph{Inverting the comparison.}
We extract a density approximation from the inverse comparison
\eqref{inp:comparison}.  Fix an integer $q\ge2$ and choose
$h_h>C_{\mathrm{cap}}+1$.  At a point in its comparison band, write
\[
 v=|y|^4H_j(y),\qquad
 t=\left(\frac{|y|^6\mu_j(y)}{k}\right)^{2/3}\ge0.
\]
Choose $0<h_l<h_h$ and $0<\xi<\min\{h_l/16,1\}$, to be decreased below.
The cap and the high implication rule out $t\ge h_h$: taking heights
increasing to $h_h$ would give $v\ge h_h-\xi>C_{\mathrm{cap}}$.
If $h_l<t<h_h$, heights approaching $t$ from below and above give
$|v-t|\le\xi$, and hence $|v_+-t|\le\xi$.
If $t\le h_l$, heights decreasing to $h_l$ in the low implication give
$v\le h_l+\xi$.  This last case includes $t=h_l$ and allows arbitrarily
negative $v$; both $t$ and $v_+$ still lie in $[0,h_l+\xi]$.
Consequently
\[
 k\bigl|t^{3/2}-(v_+)^{3/2}\bigr|
 \le
 \begin{cases}
  \dfrac32 k\sqrt{h_h}\,\xi,& h_l<t<h_h,\\[2pt]
  k(h_l+\xi)^{3/2},& t\le h_l.
 \end{cases}
\]
We may therefore choose $h_l=h_l(q)$ and $\xi=\xi(q)$ so that the
right-hand side is at most $q^{-8}$.  These choices are fixed before the
comparison's smallness threshold for $s$ is imposed.

\emph{Choosing the observation scale.}
For this fixed $q$, take $L_1=q^2$.  The admissibility condition gives
\[
 \frac{r_0}{s_\ell}
 =\eps\bigl(Z_\ell s_\ell^3\bigr)^{-1/3}\longrightarrow0.
\]
Thus, for all sufficiently large $\ell$, the largest dyadic radius not
exceeding $s_\ell/q$ has an index $j<J_\ell$ and satisfies
\[
 \frac{s_\ell}{2q}\le r_j\le\frac{s_\ell}{q}.
\]
Here $j<J_\ell$ follows from $r_{J_\ell}>s_\ell/2$ and $q\ge2$.
Moreover $4L_1r_j\ge2qs_\ell$, so the comparison covers the band
$r_j\le|y|\le qs_\ell$.
The rescaled comparison band thus extends from a radius in
$[1/(2q),1/q]$ to $q$.  On its event we have
\begin{equation}
 H_j(y)\le C_{\mathrm{cap}}|y|^{-4},\qquad
 \bigl|\mu_j(y)-k(H_j(y)_+)^{3/2}\bigr|
 \le q^{-8}|y|^{-6}
 \quad(r_j\le|y|\le qs_\ell).
 \label{scales:physical-band}
\end{equation}
The exceptional probability is at most $C_qs_\ell^{25}$.
The observations and master kernels in
\eqref{inp:observations}--\eqref{inp:packets} are fixed independently of
$L_1,h_l,h_h,\xi$.  Thus this new application of the comparison concerns
exactly the same $\mu_j,H_j$ as the fixed-parameter barrier construction.

\emph{Choosing the stages and retained data.}
Choose strictly increasing stage indices $\ell_q$ so that, for every
$\ell\ge\ell_q$, all the preceding requirements for that fixed $q$ hold,
including $C_qs_\ell^{25}\le1/q$.  On
$\ell_q\le\ell<\ell_{q+1}$ set $q_\ell=q$, and choose the deterministic
index $j_\ell$ just described.  Then $q_\ell\to\infty$,
\eqref{scales:inner-radius} holds, and the comparison events have
probability tending to one.  No bound on the growth of $C_q$, or on the
smallness thresholds as functions of $q$, is needed.

For each system, first choose the conditional-integral versions of the
fields.  At each index $j$, fix an $\mathcal F_j$-measurable full-measure
set on which all its imported identities, regularity properties, and
barrier inequalities hold.  Spatial continuity and a countable dense
family of distributional tests supply the continuum assertions; the
finitely many such sets are fixed before any data are selected.  Intersect
the comparison event at $j_\ell$ with its corresponding valid-data set
and with
\[
 \left\{\int_{\R^3}p_{j_\ell}(y)\,\dd y\le1\right\}.
\]
Call the resulting event $G_\ell$.  The intermediate error estimate
\eqref{inp:barriers} and Markov's inequality give
\[
 \Prob(G_\ell^c)
 \le\frac1{q_\ell}+C\bigl(s_\ell^{32}+s_\ell^\gamma\bigr)
 \longrightarrow0.
\]
In particular the retained-data assertion is pointwise on specified valid
versions, not an assertion left to an exceptional set after a datum is
selected.

\emph{Rescaling.}
Each conditional density has mass
$Z_\ell$, and the Coulomb kernel has degree $-1$, so
\eqref{scales:fields} gives \eqref{scales:mass-field} exactly.
The comparison \eqref{scales:physical-band} gives
\eqref{scales:band}.  Since $f(s^4u)=s^6f(u)$ for $s>0$, rescaling the
imported distributional subsolution inequality gives
\eqref{scales:invariant}, with the barriers
\eqref{scales:barriers}.  The support scales to the ball of radius
$a_\ell$, and on $G_\ell$
\[
 \int_{\R^3}p_\ell(x)\,\dd x
 =s_\ell^3\int_{\R^3}p_{j_\ell}(y)\,\dd y\le s_\ell^3,
\]
which proves \eqref{scales:error}.  The finite-system regularity follows
from the imported barrier regularity and the positive lower bound on
packet widths.  For the field remainder, a bounded integrable density has
a locally Lipschitz Coulomb potential: near the singularity the kernels
$|z|^{-1}$ and $|z|^{-2}$ are integrable in three dimensions, while away
from it both kernels are bounded.  This proves the remaining assertions.
\end{proof}
